% Protocol and provenance for the accompanying regular article.
\documentclass[11pt,a4paper]{article}
\usepackage[T1]{fontenc}
\usepackage[utf8]{inputenc}
\usepackage[english]{babel}
\usepackage[margin=25mm]{geometry}
\usepackage{amsmath,amssymb,booktabs,longtable,array}
\usepackage[numbers,sort&compress]{natbib}
\usepackage[hyphens]{url}
\usepackage[hidelinks]{hyperref}
\ifx\pdfmapfile\undefined\else
\pdfmapfile{+cm-super-t1.map}\pdfmapfile{+cm-super-ts1.map}\fi
\setlength{\emergencystretch}{2em}
\renewcommand{\thesection}{S\arabic{section}}
\renewcommand{\theequation}{S\arabic{equation}}
\renewcommand{\thetable}{S\arabic{table}}

\title{Supplementary Information\\
Stress transfer and dislocation response near a constrained
Al$_{13}$Fe$_4$ inclusion in aluminium}
\author{I.~Mikhailovskiy \and D.~Pshonkin}
\date{Moscow Polytechnic University, Moscow, Russia\\
\texttt{ilyamihailovsy@gmail.com}}

\begin{document}
\maketitle

\section{Structure and preparation}

The interface cell contains 91\,428 atoms. Matrix atoms have identifiers
1--72\,532; the remaining 18\,896 belong to the inclusion. Elemental
types are Al, Fe, Mg and Si in that order in the LAMMPS input mapping,
although the interface configurations contain no Mg or Si.
There are 87\,076 Al and 4352 Fe atoms in the complete cell.
The fcc aluminium axes are $[1\bar{1}0]$, $[11\bar{2}]$ and $[111]$,
with lattice parameter 4.05~\AA{}.
The Al$_{13}$Fe$_4$ cell is monoclinic
\cite{Liu2024CIF}; the same Jelinek interaction potential is used across
the interface and within each phase \cite{Jelinek2012}.

The initial box has $L_x=154.644253$~\AA{},
$L_y=64.482817$~\AA{} and $z$ limits $-10.0$ and
165.943041~\AA{}. The matrix construction uses 56 (111) layers.
Its upper material surface differs from the box boundary and moves
during preparation, so a single exact free-surface height is not
assumed throughout. The support layer extends to $z=20$~\AA{}.
The ridge is centred at $x=77.3221$~\AA{} and has semi-axes
35 and 20~\AA{} in the $xz$ plane.

The inclusion is scaled in-plane by factors 0.9978336 and 0.9973955
to match the periodic cell. Matrix atoms closer than 2.2~\AA{} to
inclusion atoms are removed during construction. The resulting
minimum separation is approximately 2.219~\AA{} within the cell
and 2.292~\AA{} across the periodic seam. These geometric checks
exclude initial overlaps at that distance scale; they do not
validate the relaxed interface energy or structure. The lower
region $z<6$~\AA{} is fixed.

LAMMPS uses metal units, an atomic particle style and boundaries
periodic in $x$ and $y$ and fixed in $z$
\cite{Thompson2022LAMMPS}. The neighbour skin is 2~\AA{}.
For minimisation, neighbours are rebuilt at every step.
The prepared control undergoes conjugate-gradient minimisation with
requested force two-norm 1~eV/\AA{}, followed by 6~ps at 300~K
using a Nos\'e--Hoover damping time of 0.1~ps. Langevin cooling
then lowers the target temperature to 0.5~K over 12~ps, with
damping time 0.2~ps. The velocity and cooling seeds are 4242 and
4243. Velocities are zeroed before a final minimisation whose
requested force two-norm is 0.05~eV/\AA{}. A reflecting wall at
the upper box edge is present during the thermal preparation.
This preparation defines the starting state of the paired static
calculation.

Both members of the static pair are reconstructed from that same
control snapshot. The 4816 inclusion atoms above $z=20$~\AA{}
form the displaced ridge subset. Their recorded centroid is
approximately $(77.258,32.201,28.678)$~\AA{}.
The perturbation tensor is
\begin{equation}
\mathbf E_0=
\begin{pmatrix}
0.000485&0&0.001455\\
0&-0.000970&0\\
0.001455&0&0.000485
\end{pmatrix}.
\label{eq:tensor}
\end{equation}
A reference coordinate changes by
$\delta\mathbf r=\mathbf E_0(\mathbf r-\mathbf r_c)$.
The largest imposed displacement is 0.05412~\AA{}.
The determinant of $\mathbf I+\mathbf E_0$ is
$(1+\varepsilon_0)(1-\varepsilon_0/2)^2$, with
$\varepsilon_0=0.00194$. Its departure from unity is about
$2.82\times10^{-6}$, consistent with volume conservation only
to first order in strain.

All inclusion atoms are restrained by
\begin{equation}
U_{\mathrm s}=\frac{K}{2}\sum_{i\in\mathrm{inclusion}}
 |\mathbf r_i-\mathbf r_i^{\mathrm{ref}}|^2,
\qquad K=20~\mathrm{eV}/\text{\AA}^2 .
\label{eq:spring}
\end{equation}
The support references are unshifted in both copies; the ridge
references differ by Eq.~\eqref{eq:tensor}. The spring energy is
included in minimisation. Conjugate gradients use maximum displacement
0.05~\AA{}, zero energy tolerance, force two-norm tolerance
0.02~eV/\AA{}, at most 6000 iterations and 60\,000 force
evaluations. The held control stops on force tolerance after
235 iterations, with final force two-norm 0.019779915~eV/\AA{}.
The held perturbed state also stops on force tolerance, after
310 iterations at 0.019879873~eV/\AA{}.
These values verify the stated stopping criterion.
They do not establish convergence of a stress profile with respect
to force tolerance, geometry or starting state.

\section{Stress evaluation and distances}

The stored per-atom quantities are the six components from
\texttt{compute stress/atom NULL virial}, in the order
$xx,yy,zz,xy,xz,yz$. For a control-selected set of $N$ Al atoms,
the analysis conversion is
\begin{equation}
\sigma_{\alpha\beta}
 =-\frac{1}{10N\Omega}\sum_{i=1}^{N}S_{i,\alpha\beta},
\qquad \Omega=16.6072~\text{\AA}^3,
\label{eq:stress}
\end{equation}
where the stored $S_i$ are in bar~\AA$^3$ and the result is in
MPa. This conversion makes normal compression positive and gives the
negative of the tensile-positive Cauchy tensor. Signed MD quantities in
the manuscript retain this pressure convention. The continuum comparison
uses absolute shear projections, which are unaffected by a global sign change.
No kinetic term is included. This configurational measure excludes
an interpretation of the profile as a directly measured temperature
or heat flux.

Atom identities and element types are matched between snapshots.
The control coordinates determine membership in each layer and
in the 20~\AA{} axial window. The window centre is obtained from
the midpoint of the control atom $x$ extents. Layers contain
type-1 atoms and span 4~\AA{} in $z$. A layer is retained only
when its lower edge lies above every inclusion atom in both snapshots.
Inclusion IDs exceed 72\,532 and comprise both Al and Fe.
Their maximum heights are 39.707546~\AA{} in the control and
39.718591~\AA{} in the perturbed state. The 23 retained centres run
from $r=22$ to 110~\AA{}, where $r=z-20$~\AA{}.
The last layer edges are $r=108$ and 112~\AA{}.
This choice omits layers containing inclusion aluminium and the
uppermost free-surface layers.

The tensor average is formed separately for each member of the pair
and then subtracted. With
$\mathbf D=\Delta\boldsymbol\sigma-
\operatorname{tr}(\Delta\boldsymbol\sigma)\mathbf I/3$,
the von Mises measure is $\sqrt{3\mathbf D:\mathbf D/2}$.
The largest absolute slip-system projection is
$\max_\alpha|\mathbf s_\alpha\cdot
\Delta\boldsymbol\sigma\,\mathbf n_\alpha|$.
The $xz$ projection and its sign are preserved independently.
Consequently, a small width-averaged shear can coexist with
a larger axial-window shear of the opposite sign.

The far-layer statistics are evaluated over thirteen retained
layers with $r>60$~\AA{}. The mean maximum projection is
2.477~MPa and its spatial standard deviation is 0.543~MPa.
The bins are not independent replicas of a homogeneous state.
These statistics quantify spatial variation and cannot be used
as standard errors, numerical noise estimates or a physical
detection threshold. Residual stress at the largest sampled
distance is reported explicitly.

An attenuation height can be defined only relative to a stated
stress level. If $h_m(s)$ denotes the height above the nominal
crest at which the descending axial-window magnitude reaches
$s$, linear interpolation between adjacent retained bins gives
approximately 4.69, 6.69 and 8.55~nm for $s=10$, 5 and
2~MPa, respectively. These are operational profile descriptors,
with 0.4~nm sampling, rather than independently determined
material lengths. No zero-stress attenuation height is inferred.

\section{Deformation retention and free diagnostics}

A least-squares affine map relates the positions of corresponding
Fe atoms in the two minimised configurations. After removing
translations, the fit gives a deformation gradient $\mathbf F$.
Minimum-image displacements are used along the periodic directions.
The retained strain is represented by
\begin{equation}
\mathbf E_{\mathrm{res}}=
 \tfrac12(\mathbf F^{\mathsf T}\mathbf F-\mathbf I),
\qquad
\eta=\frac{\mathbf E_{\mathrm{res}}:\mathbf E_0}
                 {\mathbf E_0:\mathbf E_0}.
\label{eq:eta}
\end{equation}
For the 1008 Fe atoms with control $z>22$~\AA{},
$\eta=0.9728$. The principal reported components are
$E_{xx}=0.000482$, $E_{yy}=-0.000887$,
$E_{zz}=0.000460$ and $E_{xz}=0.001435$.
The fit residual has root-mean-square magnitude 0.0035~\AA{}.

The formal standard error 0.0038 propagates the least-squares
coefficient covariance, using pooled residual variance over
$3N-12$ degrees of freedom. The atomic residuals are spatially
correlated; this fit statistic does not include model error or
variation between independent trajectories. A scalar projection
also does not guarantee that every tensor component retains
the same fraction of the imposed deformation.

Unrestrained minimisations are retained only as diagnostics.
For the perturbation pair without springs, the ridge projection
is $0.4364\pm0.1019$ in the saved output. Both minimisations
stop at 6000 iterations with final force two-norms
0.8660473 and 0.48161949~eV/\AA{} for control and perturbed
states, respectively, above the requested 0.02~eV/\AA{}.
A separate direct minimisation record gives
$0.2007\pm0.1064$ for a differently prepared pair.
Neither record establishes an equilibrium retained fraction,
a relaxation time or a field after magnetic removal.
Their stress profiles are excluded from the main quantitative
comparison.

\section{Loading and line tracking}

The loaded interface geometry retains 91\,428 atoms and has
the same orientation and inclusion shape as the static cell.
The initially inserted lower and upper edge lines are centred
at approximately $(139.322,49.228)$ and
$(115.922,72.611)$~\AA{} in the $xz$ plane.
Their displacements are generated by a periodic Volterra
construction. These insertion coordinates differ from the
coordinates after minimisation and initial dynamics.

The loading input first minimises the constructed cell with a
requested force two-norm of 1~eV/\AA{} and a 2000-iteration
limit. It then assigns velocities at 300~K using seed 88004
and integrates with a 1~fs step. The bottom region is fixed,
mobile atoms are coupled to a Nos\'e--Hoover thermostat with
damping time 0.1~ps, and a reflecting wall is placed at the
upper box boundary. Equal forces along $x$ on the top selected
atoms transmit the nominal shear traction. The inclusion is
either unrestrained or attached to springs of
20~eV/\AA$^2$, as specified by the case.

Let $s$ be the number of steps after the 5000-step preload,
$T_s=16\,000$ steps and $N$ the total run length. The nominal
stress in MPa is
\begin{equation}
\tau(s)=
\begin{cases}
0,&s<0,\\
q\left[\dfrac{s}{2}-\dfrac{T_s}{2\pi}
 \sin\left(\dfrac{\pi s}{T_s}\right)\right],
 &0\le s<T_s,\\
q(s-T_s/2),&s\ge T_s,
\end{cases}
\qquad
q=\frac{\tau_{\max}}{N-5000-T_s/2}.
\label{eq:ramp}
\end{equation}
The unrestrained control uses $N=101\,000$ and
$\tau_{\max}=400$~MPa. The restrained control and perturbed
runs use $N=45\,000$ and $\tau_{\max}=145.45$~MPa.
Their post-onset rates agree to the precision of these rounded
inputs. Regular snapshots are saved every 2000 steps.
The analysis tables use the actual endpoint for each run, giving
9.0909~MPa between late frames in the free control and
9.090625~MPa in the restrained pair.

DXA identifies Burgers vectors and line segments
\cite{Stukowski2012DXA,Stukowski2010OVITO}. Line positions are
unwrapped in $x$ using $L_x=154.644253$~\AA{}.
For each Burgers-vector family, a linear baseline is fitted
over steps at or above 4000 with nominal stress below 40~MPa.
The departure criterion is a residual magnitude greater than
$\max(6s_x,8~\text{\AA})$, where $s_x$ is the baseline residual
standard deviation. The subsequent two sampled positions are
checked where available; a missing line is accepted by the
stored detector. The numerical criterion therefore requires
a separate continuity check near boundaries and after
segment disappearance.

The upper-line baseline residual standard deviations are
3.8919 and 3.5046~\AA{} in the restrained control and
perturbed runs. The corresponding thresholds are
23.351 and 21.027~\AA{}. The first accepted samples are
at 38 and 40~ps, respectively. The same criterion consequently
uses different fitted baselines and displacement thresholds.
Its one-frame difference cannot be interpreted as an
independently measured 9.09~MPa strengthening increment.

The lower family disappears by 10~ps in both restrained runs.
Its later family-labelled positions can belong to seam
segments and do not provide a valid second onset for the
original line. At the last analysed control frame, 44~ps,
the total extracted length is 39.344~\AA{} in three segments.
The family-aggregated lengths are 26.191 and 13.153~\AA{};
these are not necessarily lengths of individual connected
lines. No segment is extracted from the corresponding
perturbed frame. The results distinguish absence of a
detected line from proof of an atomistic reaction pathway.

For the unrestrained control the lower-line departure is
at the 36~ps sample, 104.545~MPa, preceded by
95.455~MPa. The upper-line criterion is met at 40~ps,
122.727~MPa. The final regular analysis frame is
100~ps, 395.455~MPa; the nominal 400~MPa endpoint is
not covered by that table. Short seam signals occur
transiently around 232--241~MPa and again above
377~MPa. They do not establish persistent interface
nucleation. A scalar threshold exceeding 395.455~MPa
is not inferred from the absence of a confirmed event.

\section{Separate alloy comparator}

The alloy cell contains 92\,800 atoms with dimensions
$114.5513\times49.6022$~\AA{} in-plane and box limits
$z=-8.0$ and 283.2392~\AA{}.
Its counts are 91\,315 Al, 928 Mg and 557 Si atoms.
The nominal whole-cell concentrations are therefore
1.0~at.\% Mg and approximately 0.6~at.\% Si.
Substitution is restricted to the mobile region; the boundary
layers remain aluminium. The bottom fixed layer is
8~\AA{} thick and the top loaded layer 12~\AA{}.
The edge partners are separated by 205.768~\AA{} along
the interface normal.

The initial minimisation stopped on energy tolerance with force two-norm
10.118429~eV/\AA{} and maximum force component 1.9509711~eV/\AA{}.
It did not reach the requested force tolerance. The subsequent
300~K trajectory was completed, but this preparation does not demonstrate
mechanical equilibration of the initial dislocation arrangement.
The trajectory has
10~ps without applied shear and successive nominal
levels of 45, 55, 65 and 75~MPa, each held for 30~ps.
Each increase is smoothed over 4~ps.
The recorded net displacement of the lower line is
1.98~\AA{} from 10 to 130~ps, with a position standard deviation
of 0.63~\AA{}. The preceding zero-load interval is excluded;
the displacement across the complete 0--130~ps trajectory is 6.99~\AA{}.
The per-level net displacements
are 0.58, 1.56, $-1.53$ and 0.17~\AA{}.
The net motion over the settled portion of each level is less than
one Burgers-vector magnitude in this one solute realisation.

These observations are conditional on the sampled solute
arrangement, internal stress and short observation time.
The applied stress is not a measurement of the local
resolved traction on the dislocation. The trajectories
therefore support a statement of limited motion under
the applied programme, without establishing an intrinsic
75~MPa obstacle strength. No activation volume is
extracted from the fitted slopes.

The public comparator package contains the complete saved
atomic trajectory, including its original coordinate and stress
columns, in compressed shards under
\path{data/publication/alloy/trajectory.*.lammpstrj.gz}.
The derived 131-frame position table, summaries, initial and final
structures and log are also provided. The available initial file was
saved after the recorded run date. Its atom identities, types and fixed
boundary coordinates agree with the trajectory, but no launch-time hash
or original copy certifies its mobile coordinates. It is therefore a
candidate starting structure, not a certified exact replay input.
The saved trajectory allows independent
DXA reconstruction and reassessment of the spatial selection used
in the comparator analysis. An activation volume remains
unidentified by these observations.

The original quiet-region stress selection included Al and Si while
excluding Mg (type 2). Reanalysis of all 131 saved frames reproduces the
archived displacements and stresses. Selecting only Al or all alloy atoms
changes the mean stress on a loading level by less than 1.71~MPa.
The corresponding diagnostic estimates of $V^*/b^3$ are
$-12.32\pm21.31$ for the original selection, $-13.54\pm19.19$
for Al and $-12.32\pm19.83$ for all atoms in the pressure-sign convention.
Using tensile-positive shear reverses these fitted signs. These fits to fluctuating
positions do not determine a physical activation volume. The comparison
is recorded in \path{stageG6_stress_selection_relA.json}.

\section{Continuum and activated-glide estimates}

The sphere reference has the same small-strain tensor as
Eq.~\eqref{eq:tensor}, radius 35~\AA{}, shear modulus
26.5~GPa and Poisson ratio 0.347. The inclusion and
unbounded matrix have identical isotropic stiffness.
For the Eshelby tensor $\mathbf S$ and stiffness
$\mathbf C$, the interior stress is
\begin{equation}
\boldsymbol\sigma_{\mathrm{in}}
 =\mathbf C:(\mathbf S:\mathbf E_0-\mathbf E_0).
\label{eq:eshelby}
\end{equation}
The maximum interior slip-system projection is
41.45~MPa \cite{Eshelby1957}. This is an interior
value for a transformation-strain problem. It is not
the surface amplitude of every exterior component.
The ridge in the main calculation has different
constraints, a support and a nearby free surface,
so the comparison is not an ordering theorem or a
validation of the atomistic boundary condition.

The scalar creep calculation assumes a compact-particle
exterior amplitude $\tau_m(a/R)^3$ and volume fraction
$f$. To leading order in $f$, the spatial weight of
stress magnitudes between $s$ and $s+\mathrm ds$ is
$f\tau_m s^{-2}\mathrm ds$ for $0<s<\tau_m$.
This weight is used only under an integral for the
excess rate. It is not a normalised probability density
extended to zero stress. Equal positive and negative
stress weights and a stress-independent activation
volume give
\begin{equation}
H=f x\int_0^x\frac{\cosh u-1}{u^2}\,\mathrm du
 =f\{x\,\mathrm{Shi}(x)-[\cosh(x)-1]\},
\qquad
x=\frac{V^*\tau_m}{k_{\mathrm B}T}.
\label{eq:rate}
\end{equation}
Here $\mathrm{Shi}(x)=\int_0^x\sinh(u)/u\,\mathrm du$.
The integrand has the finite limit $1/2$ at the origin.
Particle overlap, grain-orientation distributions and
stress-dependent barrier shapes are omitted.

For the numerical sensitivity calculation,
$k_{\mathrm B}T=4.141947\times10^{-21}$~J at 300~K,
$b=2.86378246$~\AA{} and $f=0.00246$.
The volume fraction is a modelling scenario.
For a mass fraction $w=0.0035$, densities
$\rho_i=3.85$ and $\rho_m=2.70$~g/cm$^3$ give
\begin{equation}
f=\frac{w/\rho_i}{w/\rho_i+(1-w)/\rho_m}
 \simeq0.00246 .
\end{equation}
These assumed densities and composition do not establish
the volume fraction of every experimental specimen.
The activation-volume range $19b^3$--$142b^3$ is
a sensitivity interval. The reported value
$70b^3$ for an Al--Mg--Si alloy processed by
equal-channel angular extrusion provides a comparison
from a different material state \cite{Soula2022JALCOM}.

At $H=0.25$, Eq.~\eqref{eq:rate} gives
$x=6.754241794$. The required amplitudes are
62.691, 39.705, 23.823, 17.016, 11.911 and
8.388~MPa for $V^*/b^3=19$, 30, 50, 70, 100
and 142, respectively. Conversely,
$\tau_m=15.451$ and 41.45~MPa require
$V^*/b^3=77.091$ and 28.737.
The first amplitude is an axial-window ridge maximum;
the second is the sphere interior projection.
Their substitution into the scalar exterior model
is explicit and conditional. It does not calibrate
magnetostriction or demonstrate the experimental
creep rate.

\section{Records and reproducibility}

The publication data and code are located at
\par\noindent\url{https://github.com/iluuua/magnetostriction-plasticity-bounds}.
The manifest under \path{data/publication} records
the source configurations, retained columns and checksums.
It distinguishes the complete atomic-virial snapshots of
the static pair, the full-column alloy trajectory and position-only
projections of the three interface loading trajectories.
The interface projections preserve particle
identifiers, types, coordinates, frame times and cell
bounds for every saved regular frame.
These projections do not provide all original per-atom stress
and energy columns.

The stress results derive from
\path{stageG10_field_profile.json} and its associated
CSV, calculated from the held control and perturbed
atomic-virial snapshots. The layer selection,
signed projections and reported peaks can therefore
be recomputed from the primary static data.
The deformation-retention components and formal
fit error are in
\path{stageG12_eigenstrain_retention.json}.
The diagnostic free outputs are identified by
\path{stageG12_eigenstrain_retention_free.json}
and \path{stageG12_eigenstrain_retention_free_cg.json};
their lack of convergence limits their use as set
out in Section~S3.

Interface motion is recorded in
\path{stageG2_depinning_summary_G15held.json} and
\path{stageG2_depinning_summary_G15ctl_free.json}.
The associated tables
\path{stageG2_depinning_G15_ctl_G15held.csv},
\path{stageG2_depinning_G15_fld_G15held.csv} and
\path{stageG2_depinning_G15_ctl_G15ctl_free.csv}
contain frame times, nominal stresses and
family-resolved positions and lengths.
The published coordinate trajectories allow
independent reconstruction of line content.
The detector summaries require the continuity
qualification in Section~S4; later entries for a
family need not represent the same original line.

The alloy comparator uses
\path{stageG6_vstar_relA_frames.csv} and
\path{stageG6_vstar_relA_summary.json}.
Only the position history and the completed
loading range are used here. The fitted activation
volume stored in a diagnostic summary is not
treated as a physical estimate.
The sphere interior comes from
\path{stageG8_eshelby3d.json}.
The integral and inverse-stress results are
defined by Eq.~\eqref{eq:rate}, with numerical
inputs specified above; the corresponding
archived calculation is
\path{stageG5_two_scale_bridge.json}.
Narrative fields in records do not substitute
for these definitions or for convergence evidence.

The figure-generation code operates on these
configurations and records. Model renderings use
OVITO and represent simulated atomic positions.
No simulated image is presented as an experimental
micrograph. The four main figures contain the
geometry, static field, interface trajectories
and conditional activated-glide dependence.
The English labels are shared by both language
versions.

\bibliographystyle{unsrtnat}
\bibliography{references}
\end{document}
